\begin{helper}[Last-index component-intersection shortening]
Let \((\varphi,T)\) be a minimum bad pair with
\(\rho(T)=|E(T)|\).  Let \(S\) be an exact initial segment of \(T\) with
\(|E(S)|+2=|E(T)|\).  Suppose distinct colors \(\alpha,\beta\) are
missing on \(S\), \(\beta\) is also missing at the last vertex \(y_n\)
of \(T\), and \(y_n\notin V(S)\).  Then the
\((\alpha,\beta)\)-component rooted at \(y_n\) cannot meet \(V(S)\).
\end{helper}
\begin{proof}
Assume that the component meets \(V(S)\).  Unfolding
\noderef{TashTreeMinimumBadPair} and \noderef{TashTreeBadPair}, obtain
the proper \(k\)-triple for \(\varphi\), the Tashkinov condition on
\(T\), and the universal exclusion of a smaller bad pair.  Restrict the
Tashkinov condition to \(S\): for every formal edge index of \(S\) at
least one, \noderef{TreeSequenceInitialSegment} identifies its exact
edge and every earlier vertex with the corresponding entries of \(T\),
so the witness furnished by \noderef{TashkinovTree} for \(T\) is also a
witness in \(S\).

Apply \noderef{TwoColorComponentTashkinovSplice} to \(S\), the two
tree-missing colors, the outside vertex \(y_n\), and the assumed
intersection.  It returns a vertex \(u\in V(S)\) and a tree-sequence
\(S'\) which properly extends \(S\), ends at \(y_n\), and is connected
to \(S\) by a \noderef{TashkinovAugmentationSequence}.  The splice
statement records every appended exact edge-node and vertex-node in the
chosen component, and says that the first appended edge joins \(u\) to
its new vertex while every later appended edge joins consecutive new
vertices.  Since \(S\) is Tashkinov, applying
\noderef{TashkinovAugmentationSequencePreservesTashkinovTree} shows that
\(S'\) is Tashkinov under the unchanged coloring.

The hypothesis \(\beta\in\noderef{TreeMissingColors}(\varphi,S)\)
gives a vertex \(v\in V(S)\) missing \(\beta\).  The initial-segment
clause of the splice puts \(v\) in \(V(S')\), while its last-vertex
clause puts \(y_n\) in \(V(S')\).  These vertices are distinct because
\(y_n\notin V(S)\), and \(y_n\) also misses \(\beta\).  Thus
\noderef{ColoredVertexSetElementary} fails on
\noderef{TreeSequenceVertexSet} of \(S'\).  The original \(k\)-triple,
the Tashkinov condition just proved, and this repeated missing color give
\noderef{TashTreeBadPair} for \(\varphi,S'\).

Set \(q=|E(S)|+1\).  It is positive, and the splice grows \(S\) by at
least one edge, so \(q\le |E(S')|\).  For every formal edge index
\(r\ge q\), one has \(r>|E(S)|\); hence it is not the first appended
edge, and the final endpoint equation in the splice statement identifies
its ends with the consecutive displayed vertices of \(S'\).  These are
exactly the clauses of \noderef{TashTreeSuffixPath} at \(q\).  By the
leastness in \noderef{TashTreePathIndex},
\[
 \rho(S')\le |E(S)|+1=|E(T)|-1<|E(T)|=\rho(T).
\]
Thus the first coordinate of \noderef{TashTreeRankLess} strictly
decreases.  The smaller bad pair contradicts the universal clause of
\noderef{TashTreeMinimumBadPair}, proving that the component cannot meet
\(V(S)\).
\end{proof}
